Akustiske undersøkingar støttar bestandsestimering i fiskeri og kartlegging av leveområde, der botnregistrering skil tilbakespreiing frå vassøyla frå botnekkoet \cite{kenny_overview_2003,ona_mitson_1996,thomas_trawl_2024}. Operative detektorar kan feile utan varsel når føresetnadene deira ikkje held, slik at feil blir førte vidare til analysar som behandlar den registrerte botnen som referanse. Som ein konseptstudie undersøker denne oppgåva om djup læring kan gjenskape algoritmiske referanseetikettar i Simrad EK80-prosessering og kvantifisere si eiga pålitelegheit med konformal prediksjon.

Vi evaluerer U-Net-modellar for botnsegmentering på tofrekvente Simrad EK80-vidbandsekkogram med kalibrert volumtilbakespreiingsstyrke ($S_v$) ved 38\,kHz og 200\,kHz, saman med $S_v^{38}-S_v^{200}$. Etikettane blir genererte på 38\,kHz-kanalen ved å ta gjennomsnittet av to tradisjonelle detektorar og jamne resultatet over tid. Dei er uverifiserte akustiske estimat, ikkje målte etikettar. Dataa kjem frå M/F Bast\o{} Electric, ei ferje som kryssar Oslofjorden om lag 24 gonger dagleg, og gjev gjentekne opptak av ei rute under skiftande vêr, tidevatn og biologiske tilhøve.

På testperioden frå oktober 2024, som omfatta 378 filer som ikkje inngjekk i treninga, men som kryssa den same botnen, nådde Bast\o{} Standard U-Net ein gjennomsnittleg snitt-over-union (IoU) på 0,9972, ein median på 0,998 og eit gjennomsnittleg djupneavvik på 0,57\,m. Innanfor dei respektive evalueringsprotokollane gav større enkodarar, ImageNet-førtrening, ensemble over tilfeldige frø og endringar i tapsfunksjonen lite evidens for avgjerande forbetringar. I ein separat ablasjon på 100 filer betra augmenteringar som etterliknar feil i datainnsamlinga djupnestabiliteten innanfor dette deldatasettet.

Konformal prediksjon gav kalibrerte band innanfor det same geografiske regimet, men ingen generell garanti for pålitelegheit. Eit supervisert split-konformalt band kalibrert på eit prediksjonsmellomlager frå januar 2025 dekte dei tilsvarande testmellomlagera frå februar og juni 2025 konservativt, med 92,6\% dekning mot eit mål på 90\%. Eit etikettfritt band basert på gjentekne kryssingar identifiserte upålitelege strekningar utan manuelle etikettar, heldt dekninga på 90,2\% og reduserte median halvbreidd frå 7,59 til 3,59\,m. Ved overføring til ei NOAA-undersøking av stillehavslysing vart både segmenteringssamsvaret og dekninga med Bast\o{}-kalibrering svakare. Ny kalibrering på 20 lysingfiler gav om lag 90\% dekning, men breiare band, medan finjustering på lysingdata gjenoppretta både samsvar og dekning som domenespesifikk tilpassing, ikkje som ein universell modell. Bidraget er ein kasusstudie av lært EK80-botnsegmentering med kalibrert konformal diagnostikk innanfor ei fast rute, avgrensa av etikettkvalitet og avgrensa datasettmangfald.
